% paper/sections/08_2_fccd_bf.tex
%
% §8.5  FCCD 面フラックス部分系の構成
%
% ═══════════════════════════════════════════════════════════════════════════════
\subsection{FCCD 面フラックス部分系の構成}
\label{sec:fccd_bf_sub_system}
% ═══════════════════════════════════════════════════════════════════════════════

§~\ref{sec:fccd_def} で導出した FCCD 面フラックス表現
$\FaceJet{q}=(q_f,q'_f,q''_f)$ は，BF 部分系の
P-1..P-4 を\textbf{同時に}満たす面表現である．
圧力勾配，速度補正，HFE/GFM ジャンプ行，毛管 cochain を同じ面フットプリントへ写すことで，
連続平衡
$-\bnabla p+\mathbf{f}_\sigma=\mathbf{0}$
を離散空間でも同じ面釣合として扱える．

% -------------------------------------------------------------------------------
\phantomsection
\label{sec:fccd_face_residual_order}

本稿の構成では {\FCCD} 面フラックス表現 $\FaceJet{q}=(q_f,q'_f,q''_f)$ が
圧力勾配・速度補正・HFE 片側データ・圧力ジャンプ条件の 4 経路で共有される（(R1)-(R4)）．
これにより面フットプリントは 4 経路で一致し，
\begin{equation}
  \mathcal{G}p := D^\FCCD p,
  \qquad c_f(j_{gl}) := A_f B_f(j_{gl})
  \qquad\Longrightarrow\qquad
  a_f(p;c)=A_f\,\mathcal{G}p-c_f
  \label{eq:bf_fccd_unification}
\end{equation}
となる．Balanced--Force の判定対象は，個別の節点勾配ではなく
同じ face complex 上の実効補正 cochain $a_f$ である．
標準 pressure-jump 閉包では $c_f$ を
§~\ref{sec:split_ppe_capillary_range_projection} の
surface-energy virtual-work/Riesz 構成で定め，
同じ面複体の Hodge gate で静的平衡か非平衡駆動かを判定する．
圧力勾配値域への射影はこの gate の診断代表であり，
本番補正前に divergence-free 毛管 cochain を一律に消すための操作ではない．

% -------------------------------------------------------------------------------
\subsubsection{\texorpdfstring{面フラックス表現 $\FaceJet{p}$ による BF 構築}{面フラックス表現 FaceJet(p) による BF 構築}}
\label{sec:face_jet_bf_construction}
% -------------------------------------------------------------------------------

面フラックス表現 $\FaceJet{p}=(p_f,p'_f,p''_f)$ は
P-1..P-4 を同一の面表現で具現する．
変密度 PPE の\textbf{面フラックス}は
\begin{equation}
  F^p_f := A_f\,\FaceJet{p}_1 = A_f\,p'_f,
  \label{eq:fccd_bf_face_flux}
\end{equation}
ここで $A_f$ は逆密度，切断面係数，局所距離，メトリクスを含む
P-4 の単一面係数である（一様格子では $\betaf=(1/\rho)_f$ に縮約）．
PPE は FVM セル積分で
\begin{equation}
  F^p_{i+1/2} - F^p_{i-1/2}
    = \int_{x_{i-1/2}}^{x_{i+1/2}}\! S\,dx,
  \qquad S = \frac{1}{\Delta t_{\mathrm{proj}}}\,D_h^\text{bf}\,\bu^*,
  \label{eq:fccd_bf_cell_balance}
\end{equation}
として組み立て，速度補正は
\begin{equation}
  u_f^{n+1}
    = u_f^* - \Delta t_{\mathrm{proj}}\,A_f\,p'_f
    \quad\text{（同じ }p'_f\text{）．}
  \label{eq:fccd_bf_corrector}
\end{equation}
\textbf{同じ $p'_f$} が PPE 組立と補正で共有されるため P-2 が構造的に成立．

% -------------------------------------------------------------------------------
\subsubsection{界面ジャンプ GFM 補正（標準 pressure-jump 閉包）}
\label{sec:fccd_bf_gfm_integration}
% -------------------------------------------------------------------------------

界面 $\Gamma$ を跨ぐ面 $f$ で面フラックス表現を
ジャンプ条件 $j_{gl}=p_g-p_l=-\sigma\kappa_{lg}$，$[A_f\partial_n p]_\Gamma=0$
（非相変化）で補正する．標準経路では，ジャンプ演算子
$B_f(j_{gl})$ を同じ局所物理距離 $H_f$ と同じ面係数 $A_f$ から定義し，
\begin{equation}
  a_f(p;c)
    = A_f G_f p - A_f B_f(j_{gl})
  \label{eq:fccd_bf_gfm_face}
\end{equation}
を補正段階の実効 face cochain とする．
相内の滑らかな離散化は保ち，ジャンプ項だけを面フラックスへ加える．
HFE はこの GFM/affine-jump 行に必要な片側 Hermite データを供給し，
$c_f=A_f B_f(j_{gl})$ または表面エネルギーから得る
Riesz cochain は同じ face slot で速度補正へ渡す．
本稿の標準構成では GFM/affine-jump + HFE + DC 残差契約 +
変分毛管 Hodge gate を用いる．
固定回数 $k=3$ は §~\ref{sec:defect_correction_ppe} の成分検証での実用上限であり，
補正段階の標準契約は高次残差が規定許容値へ落ちることである．

% -------------------------------------------------------------------------------
\subsubsection{静水圧テストによる BF 精度検証}
\label{sec:fccd_bf_hydrostatic}
% -------------------------------------------------------------------------------

\noindent\textbf{検証問題}：重力 $\mathbf{g}=-g\hat{\mathbf{z}}$ 下の静止 2 層流体．
解析解は $p_\text{stat}(z) = -\rho(\phi(z))\,g\,z$，$\bu=\bm{0}$．

\noindent\textbf{BF 残差測定}：
重力加速度を $\bm{a}^g_f$，静水圧の面勾配を
$(G_h^\text{bf}p_\text{stat})_f$ として，
一ステップ補正で速度へ注入される離散加速度残差を測る：
\begin{equation}
  \varepsilon_\text{BF}
    := \Delta t_{\mathrm{proj}}
       \max_f\,
       \Bigl|\,
         \bm{a}^g_f
         - A_f\,(G_h^\text{bf}p_\text{stat})_f
       \Bigr|.
  \label{eq:fccd_bf_hydrostatic_metric}
\end{equation}
これは $\bu^*=\bm{0}$ から開始し，重力加速度と
圧力補正を同じ面スロットで同時に適用した
1 ステップ釣合確認後の $\max_f|u_f^{n+1}|$ と同値である
（§~\ref{sec:time_buoyancy} の
balanced-buoyancy 分解と同じ考え方）．
理想離散 BF では $\varepsilon_\text{BF}=0$．このテストの読み方は
数値値そのものではなく，P-1/P-2/P-4 のいずれを破ったかを
一ステップ速度残差として分離できる点にある：
\begin{center}
\begin{tabular}{@{}p{0.30\textwidth}p{0.45\textwidth}p{0.18\textwidth}@{}}
\hline
構成 & 検出する違反 & 期待される判定 \\
\hline
非整合な節点勾配 + 事後面平均 & 面ロカス不一致，補正 $G_h$ 不一致 & FAIL \\
同一 face-adjoint pair & PPE と補正の面複体共有 & PASS \\
pressure-jump Hodge gate & 静的平衡でない毛管 Hodge 成分の消去 & FAIL \\
\hline
\end{tabular}
\end{center}
\noindent pressure-jump 経路の実測値は §~\ref{sec:bf_static_droplet} の
face-balance gate で示す．そこでは raw 毛管 cochain の非勾配成分と，
静的平衡で消えるべき成分・非平衡界面の物理駆動成分を分けて評価する．

\noindent\textbf{重力 BF 検証}：
本表は重力平衡での BF 残差を測定する
（CSF BF の静止液滴テスト §~\ref{sec:bf_static_droplet} とは独立；
CSF 側は $f_\sigma$ と $G_h p_\sigma$ の相殺が主論点）．
この診断の検証は §~\ref{sec:verify_split_ppe_dc} と §~\ref{sec:bf_static_droplet} に分離する．

詳細検証は §~\ref{sec:verify_split_ppe_dc}（要素検証）と
§~\ref{sec:bf_static_droplet}（静止液滴 gate）．

% -------------------------------------------------------------------------------
\subsubsection{面演算子閉包 A/B/C の位置付け}
\label{sec:fccd_bf_strategy_table}
% -------------------------------------------------------------------------------

\begin{center}
\begin{tabular}{@{}p{0.20\textwidth}p{0.40\textwidth}p{0.30\textwidth}@{}}
\hline
構成 & 内容 & 用途 \\
\hline
\textbf{A：標準 affine-jump face complex} &
$D_fG_A$，$B_f(j_{gl})$，変分毛管 Riesz cochain，
および component saddle / Hodge gate を同じ pressure-adjoint face complex で構築．
HFE は片側 Hermite データを供給．
& 高密度比・低 We の標準 pressure-jump 閉包．
\\
\textbf{B：CSF/BF 縮約} &
CSF 体積力を予測段へ入れ，PPE + 補正 + $f_\sigma$ を同じ面層に置く．
& 低密度比・一括 PPE の縮約検証経路．
\\
\textbf{C：FCCD DNS 極限} &
BF 面整合を CLS，運動量，分相 PPE，HFE/GFM へ広げる．
& 標準経路の高次極限としての研究構成．
\\
\hline
\end{tabular}
\end{center}

\noindent\textbf{本稿の標準構成}：
バルクの高次 CCD 系演算子で移流・拡散精度を確保しつつ，
PPE，速度補正，HFE/GFM ジャンプ行，毛管 face cochain は同じ BF 面演算子へ集約する．
標準 pressure-jump 閉包では affine-jump face law と変分毛管 Hodge gate を用い，
CSF/BF は低密度比・一括 PPE の縮約診断として扱う．
純 FCCD DNS はこの面整合を全演算子へ広げた高次極限として位置付ける
（§~\ref{sec:pure_fccd_dns}）．

% -------------------------------------------------------------------------------
\subsubsection{BF アンチパターン一覧}
\label{sec:fccd_bf_anti_patterns}
% -------------------------------------------------------------------------------

\begin{tcolorbox}[warnbox,title={避けるべき BF アンチパターン}]
\begin{itemize}[leftmargin=1.5em,noitemsep]
  \item 高次 $G_h p$，低次毛管 cochain（切断誤差非対称）
  \item 節点 CCD 全量 $\kappa_{lg}$（界面帯ノイズ増幅）
  \item PPE $G_h$ と補正 $G_h$ が別演算子（F-2）
  \item 毛管 cochain が節点，$G_h p$ が面（F-1）
  \item $A_f/\betaf$ 混用（PPE 調和 / 補正算術 / ジャンプ cochain 別メトリクス；F-4）
  \item 界面越え CCD ステンシル無補正（F-5）
  \item 表面エネルギーの仮想仕事を経ない raw 曲率 cochain を物理力として採用する
  \item 非平衡界面の毛管 Hodge 成分を range projection で本番補正前に消す
\end{itemize}
\end{tcolorbox}

% ═══════════════════════════════════════════════════════════════════════════════
